% SPDX-License-Identifier: Apache-2.0
%
% TxHDL from zero: a blinky in Bazel, from an empty directory to a
% netlist, for somebody who has neither the tree nor Bazel. Built by
% //docs:zero. Every file shown is included from tutorial/blinky/,
% whose design the tree builds against its own library as a check,
% and every printout is what that build produced by running it. The
% pin in its MODULE.bazel is checked by building the workspace outside
% the tree, which no action here does (issue 733).
\documentclass[journal,9pt]{IEEEtran}

\newcommand{\listingsize}{\fontsize{6}{7}\selectfont}
% SPDX-License-Identifier: Apache-2.0
%
% The house preamble, shared by every document in this directory. Loaded
% after \documentclass, before \title. Keeping it in one file is what
% stops two articles drifting apart in font, listing style or figure
% style.
% Computer Modern. IEEEtran selects Times (ptm) by default, so the face has
% to be chosen explicitly.
%
% Latin Modern is the Computer Modern variety used here, and the choice is
% forced rather than aesthetic. Under T1 encoding, plain `cmr` has no Type 1
% outlines unless cm-super is installed, and it is not in the pinned
% distribution, so pdflatex falls back to EC bitmaps and every glyph in the
% PDF becomes a Type 3 font. That was measured: the first build produced a
% document with 10 Type 3 fonts and no Type 1. Latin Modern is Computer
% Modern redrawn with a full T1 repertoire and real .pfb outlines, so the
% same design comes out scalable.
\usepackage[T1]{fontenc}
\usepackage[utf8]{inputenc}
\usepackage{lmodern}
% microtype is not in the pinned TeX distribution. emergencystretch alone
% keeps the two-column measure from producing overfull lines.
\setlength{\emergencystretch}{3em}

\usepackage{amsmath}
\usepackage{amssymb}
\usepackage{amsthm}
\usepackage{booktabs}
\usepackage{array}
% enumitem is not in the pinned TeX distribution; the list spacing below
% does what \setlist would have done.
\usepackage{float}
\usepackage{xcolor}
\usepackage{listings}
\usepackage{tikz}
\usetikzlibrary{arrows.meta,positioning,fit,backgrounds,calc}
\usepackage[hidelinks,bookmarks=true]{hyperref}

% Numbered, definition-style box for a rule the reader is meant to apply.
\theoremstyle{definition}
\newtheorem{rulebox}{Rule}
\newtheorem{example}{Example}

% Unnumbered asides.
\theoremstyle{remark}
\newtheorem*{tip}{Tip}
\newtheorem*{pitfall}{Pitfall}

% Tighter lists, in place of enumitem's \setlist.
\setlength{\topsep}{2pt}
\setlength{\itemsep}{1pt}
\setlength{\parsep}{0pt}

% Code in the running text. The typewriter face under T1 ligates `<<`
% and `>>` into guillemets, so a shift printed as \code{<<} came out as
% a French quotation mark (issue 571). microtype's \DisableLigatures is
% not in the pinned distribution, but the primitive it rests on is
% pdfTeX's own: \pdfnoligatures strips every ligature from the font it
% is given, and the current font inside \texttt is the typewriter face
% at the size in use, so each size is stripped the first time \code
% sets it. Code wants no ligature at all: two quotes are two quotes.
\newcommand{\code}[1]{\texttt{\ifdefined\pdfnoligatures\pdfnoligatures\font\fi #1}}
\newcommand{\term}[1]{\emph{#1}}

% Syntax highlighting. Four roles, distinguishable in colour and still
% distinguishable in greyscale, because an IEEE article gets printed: the
% keyword is the only bold role, the comment is the only italic one, and the
% two remaining roles differ in lightness as well as in hue.
\definecolor{kwblue}{RGB}{20,60,150}
\definecolor{cmtgreen}{RGB}{90,120,90}
\definecolor{strred}{RGB}{150,50,40}
\definecolor{numgray}{gray}{0.45}
\definecolor{framegray}{gray}{0.70}
\definecolor{bgshade}{gray}{0.97}
% A darker fill for figure cells. The listing background has to stay very
% light to keep code readable; a figure cell has to be clearly shaded or the
% distinction it draws is invisible in print.
\definecolor{cellshade}{gray}{0.90}

% The listing size is the document's choice, because it depends on the
% column: a two-column article sets `\listingsize` to 6pt and keeps its
% listings within 70 columns, or to 5.6pt when it includes source that
% rustfmt sets at 80, which is what fits a column's frame (issue 1061);
% a one-column document keeps the body size and includes source files
% at 80. `//docs:frame_test` checks every listing line against its frame. Lines are never broken; a line that does not fit
% overflows the frame, which is visible, rather than wrapping, which is
% not.
\providecommand{\listingsize}{\scriptsize}

% `'` is deliberately not a string delimiter: the language uses it for loop
% labels, as Rust does, so treating it as a quote would swallow the rest of
% the line.
\lstdefinestyle{txhdl}{
  basicstyle=\ttfamily\listingsize,
  keywordstyle=\bfseries\color{kwblue},
  commentstyle=\itshape\color{cmtgreen},
  stringstyle=\color{strred},
  numberstyle=\tiny\color{numgray},
  morekeywords={mod,use,pub,const,type,struct,enum,trait,impl,fn,async,let,mut,
    match,if,else,loop,while,for,in,break,continue,return,as,await,
    module,bus,channel,transaction,protocol,pipeline,flow,seq,config,
    port,stage,pipe,instance,connect,bind,tag,domain,blackbox,foreign,
    property,role,signal,handler,true,false,
    sequence,interface,modport,implements,package,import,var,wire,func,proc,
    case,default,next,fork,branch,join,state,fsm},
  morecomment=[l]{//},
  morestring=[b]",
  showstringspaces=false,
  columns=fullflexible,
  keepspaces=true,
  breaklines=false,
  % Framed, so a listing is bounded rather than floating in the column.
  frame=single,
  rulecolor=\color{framegray},
  framesep=3pt,
  backgroundcolor=\color{bgshade},
  xleftmargin=3pt,
  xrightmargin=3pt,
  aboveskip=6pt,
  belowskip=6pt,
  % Every listing is numbered and captioned, so it can be referred to by
  % number. `caption` without `float` numbers the listing and leaves it
  % where it was written, which is what a listing in running prose needs.
  captionpos=b,
  abovecaptionskip=2pt,
  belowcaptionskip=6pt,
}

% Figures drawn as text. Framed the same way, and with the highlighting
% switched off, because a box-and-arrow diagram is not code and half its
% words turning blue reads as an accident. Setting a style adds keys rather
% than clearing the ones already set, so the three roles are neutralised by
% hand rather than by leaving them out. Program output is printed in this
% style too, and a netlist's assignment can run past the frame, so a line
% that does not fit is broken and the rest indented; source never is.
\lstdefinestyle{ascii}{
  basicstyle=\ttfamily\listingsize,
  keywordstyle=\ttfamily,
  commentstyle=\ttfamily,
  stringstyle=\ttfamily,
  columns=fullflexible,
  keepspaces=true,
  breaklines=true,
  breakindent=2em,
  frame=single,
  rulecolor=\color{framegray},
  framesep=3pt,
  backgroundcolor=\color{bgshade},
  xleftmargin=3pt,
  xrightmargin=3pt,
  aboveskip=6pt,
  belowskip=6pt,
}
\lstset{style=txhdl}

% House TikZ styles. `shorten` lives on the arrow styles rather than on each
% arrow, so every arrow in the document gets the gap, including ones added
% later. TikZ clips a path at a node's border, so without it an arrowhead
% butts against the box with no gap at all. Keep `node distance` well above
% twice the shorten, or the arrow shrinks to nothing.
\tikzset{
  box/.style={draw,rounded corners=2pt,align=center,inner sep=4pt,
              font=\footnotesize},
  boxf/.style={box,fill=bgshade},
  lbl/.style={font=\scriptsize\itshape,align=center},
  fl/.style={-{Stealth[length=4pt]},shorten >=2.5pt,shorten <=2.5pt},
  fld/.style={fl,dashed},
}


% The contents list: the class gives a section's number 2.75em, which
% a Roman numeral past thirty overruns onto the title, so the box is
% wider here, and a subsection's entry starts where a title does.
\makeatletter
\def\l@section#1#2{\addpenalty{\@secpenalty}\addvspace{1.0em plus 1pt}%
    \@tempdima 5.75em \begingroup \parindent \z@ \rightskip \@pnumwidth%
    \parfillskip-\@pnumwidth {\bfseries\leavevmode #1}\hfil\hbox to\@pnumwidth{\hss #2}\par%
    \endgroup}
\def\l@subsection{\@dottedtocline{2}{5.75em}{5em}}
\def\l@subsubsection{\@dottedtocline{3}{10.75em}{5.5em}}
\makeatother


\InputIfFileExists{buildstamp.tex}{}{}
\providecommand{\buildstamp}{Draft build (outside the Bazel flow).}

\title{TxHDL from Zero: A Blinky in Bazel}

\author{Filip Filmar and Dragiša Janković%
\thanks{Both authors are independent. Filip Filmar is the
corresponding author, at f@hdlfactory.com; Dragiša Janković is
reached at linkedin.com/in/dragisa-jankovic.}%
\thanks{This bootstrap guide configures a standalone hardware project
from an empty workspace: with only Bazelisk and Git installed and without
cloning the monorepo, a developer compiles, simulates, and synthesizes an LED
blinky to Verilog. Listings are extracted directly from \code{tutorial/blinky/}
in the repository \code{HDL/txhdl}, and execution logs reflect live build
output. Architectural foundations appear in~\cite{tutorial} and the cheat
sheet~\cite{cheatsheet}; the project index~\cite{cover} catalogues all companion
reports.}%
\thanks{The authors used a large language model, Claude, as an
assistant in exploring concepts and in drafting and constructing
the documents and implementations; every repository commit documents
this collaboration and includes the prompt sequence verbatim.}%
\thanks{\buildstamp}}

\markboth{TxHDL from Zero}{Filmar and Janković: TxHDL from Zero}

\begin{document}
\maketitle

\begin{abstract}
TxHDL is a hardware description language implemented directly as a Rust
library, built and lowered through Bazel. This guide starts from an empty
directory and delivers an operational blinky module that simulates in software,
renders an execution trace, and derives synthesizable Verilog as a build
artifact. The setup requires exactly five configuration and design files,
predicated on two host utilities: Bazelisk and Git.
\end{abstract}

\section{What You Need}
\label{sec:z-need}

Two prerequisites suffice: Bazelisk, provisioned as \code{bazel} to fetch the
pinned Bazel release automatically, and Git. All compilers, linkers, Rust
toolchains, and runtime dependencies are resolved hermetically by Bazel. No
external Verilog toolchain is required; synthesizable RTL is derived and printed
directly by the compiled binary.

\section{The Workspace}
\label{sec:z-files}

A standalone workspace requires five files: four environment configuration
files detailed below, and the design file discussed in
Section~\ref{sec:z-blinky}. First, \code{.bazelversion} pins the Bazel release
guaranteeing reproducible execution:

\lstinputlisting[style=ascii,caption={\code{.bazelversion}: one line.},label={lst:z-version}]{tutorial/blinky/.bazelversion}

Second, \code{.bazelrc} declares registry endpoints for external module
resolution, prioritizing the project registry over the Bazel Central Registry
(BCR):

\lstinputlisting[style=ascii,caption={\code{.bazelrc}: where modules come from.},label={lst:z-rc}]{tutorial/blinky/.bazelrc}

Third, \code{MODULE.bazel} defines the workspace, binding external dependencies
on \code{rules\_rust} and the TxHDL GitHub mirror pinned to an immutable commit.
Because Bzlmod toolchain registration is governed by the root module, the
workspace explicitly configures Rust 1.90.0 to match TxHDL's compiler baseline:

\lstinputlisting[style=ascii,caption={\code{MODULE.bazel}: two dependencies, one of them by commit, and the Rust compiler.},label={lst:z-module}]{tutorial/blinky/MODULE.bazel}

Fourth, \code{BUILD.bazel} declares build targets. The primary target is a
\code{rust\_binary} linking TxHDL runtime libraries and procedural macros; a
companion \code{genrule} invokes the binary with \code{--verilog} to capture
synthesized RTL into \code{blinky.v}:

\lstinputlisting[style=ascii,caption={\code{BUILD.bazel}: the binary, its netlist, and the sources.},label={lst:z-build}]{tutorial/blinky/BUILD.bazel}

\section{The Blinky}
\label{sec:z-blinky}

The hardware design comprises a 32-bit counter register and a scalar output.
Sequential state resides in struct fields, while temporal behavior is
implemented within \code{Unit::run}: an infinite loop initiated by an edge
wait, making each iteration represent one clock cycle. Macro \code{when!}
handles counter wrap-around, and combinational output \code{led} is driven via
\code{set}. The procedural attribute \code{\#[lower]} derives synthesizable
Verilog via \code{Blinky::verilog}. Embedded function \code{main} executes the
simulation across two full periods, printing character output or emitting
Verilog on request.

\lstinputlisting[caption={\code{blinky.rs}: the design, and a \code{main} that runs it or prints it.},label={lst:z-rs}]{tutorial/blinky/blinky.rs}

\section{Build and Run}
\label{sec:z-run}

Execute the design directly from the workspace root:

\begin{lstlisting}[style=ascii]
bazel run //:blinky
\end{lstlisting}

The initial run fetches Bazel, the Rust toolchain, and TxHDL with all
transitive dependencies, compiling all artifacts in approximately two minutes.
Subsequent runs execute in seconds from cache. Simulation prints one character
per cycle, asserting the LED for the initial eight cycles of each sixteen-cycle
period:

\lstinputlisting[style=ascii,caption={What \code{bazel run //:blinky} printed when this document was built.},label={lst:z-out}]{out_zero.txt}

\section{The Netlist as a File}
\label{sec:z-netlist}

Synthesize standalone Verilog via:

\begin{lstlisting}[style=ascii]
bazel build //:blinky.v
cat bazel-bin/blinky.v
\end{lstlisting}

The emitted module contains the registered counter, synchronous reset and wrap
guards, and continuous assignment driving the output, with constants folded
at compile time:

\lstinputlisting[style=ascii,caption={\code{blinky.v}, as the build wrote it.},label={lst:z-v}]{tutorial/blinky/blinky.v}

\section{From Here}
\label{sec:z-next}

Modifying \code{PERIOD} updates both simulation traces and synthesized RTL
immediately. Deriving \code{Trace} on additional registers exposes them in
waveforms, as detailed in the step-by-step tutorial~\cite{tutorial}. Companion
references include the one-page cheat sheet~\cite{cheatsheet}, verified
examples~\cite{examples}, and the processor implementation~\cite{vreteno}.
Updating the pinned commit in \code{MODULE.bazel} pulls upstream enhancements
seamlessly.

\begin{thebibliography}{9}

\bibitem{tutorial}
F.~Filmar and D.~Janković, ``TxHDL, step by step,'' project
repository \code{HDL/txhdl}, \code{//docs:tutorial}, 2026.

\bibitem{cheatsheet}
F.~Filmar and D.~Janković, ``TxHDL cheat sheet,'' project repository
\code{HDL/txhdl}, \code{//docs:cheatsheet}, 2026.

\bibitem{examples}
F.~Filmar and D.~Janković, ``TxHDL by example,'' project repository
\code{HDL/txhdl}, \code{//docs:examples}, 2026.

\bibitem{vreteno}
F.~Filmar and D.~Janković, ``Vreteno: an RV32IMAC core in TxHDL,''
project repository \code{HDL/txhdl}, \code{//docs:vreteno}, 2026.

\bibitem{cover}
F.~Filmar and D.~Janković, ``TxHDL documents,'' project repository
\code{HDL/txhdl}, \code{//docs:cover}, 2026.

\end{thebibliography}

\end{document}
